% GENERATED FILE. Edit canonical lessons, then run npm run generate:books.
% canonical-source-hash: fnv1a64:427ca3353d8a9fea
% canonical-lessons: BN-C25-ke-suffix, BN-C25-amake, BN-C25-tomake, BN-C25-ke-who, BN-R25-the-one-it-happens-to

\chapter{The One It Happens To}
\label{ch:bn-the-one-it-happens-to}
\hlchaptermodality{36}

\begin{chapteropening}
\textbf{By the end of this chapter:} \emph{I can mark the person a thing is done to, say me, you and him as objects, and ask who?}

The object ending goes on living nouns only — the animacy question paying for the third time — and fills the slot in chapter nine's dative frame, which had been taught with nothing to put in it for twenty-seven chapters.
\end{chapteropening}

\section[-ke]{\bn{কে} (-ke) — the one it happens to}

\label{lesson:BN-C25-ke-suffix}



\begin{quote}
Which group is \textbf{\bn{ভাই}} in and which is \textbf{\bn{জামা}} in? Then say \textbf{our family}.

\begin{itemize}
\raggedright
  \item \emph{Recall:} two chapters back, the word for \emph{he} on the careful rung — \emph{tini}
\end{itemize}
\end{quote}

\subsection*{You'll want to know: \bn{কে}}
\begin{quote}
\textbf{\bn{ভাইকে}} — \emph{bhāike} — \textbf{to the brother}, or \textbf{the brother}, as the one it is done to.
\end{quote}

\begin{quote}
\textbf{\bn{বন্ধুকে}} — \emph{bôndhuke} — \textbf{the friend}, as the one it is done to.
\end{quote}

Another ending on the back of the noun, and it marks something new: that this person is not doing the action but \textbf{receiving} it. English marks that with word order alone; Bengali marks it on the word.

And it is choosy in a way you can now predict. \textbf{\bn{কে}} goes on \textbf{\bn{জীব}} nouns — people — and not on things:

\begin{itemize}
\raggedright
  \item \textbf{\bn{ভাইকে}} — yes
  \item \textbf{\bn{জামাকে}} — no; a shirt stays \textbf{\bn{জামা}}
\end{itemize}

The animacy split you learned three chapters ago is doing this too. That one question keeps paying.

\begin{grammarlens}[title={Grammar Lens: the same slot, still one at a time}]
Line the endings up and the system is visible:

\noindent
\begin{tabularx}{\linewidth}{@{}>{\raggedright\arraybackslash}X>{\raggedright\arraybackslash}X>{\raggedright\arraybackslash}X@{}}
\toprule
\textbf{ending} & \textbf{what it says} & \textbf{goes on} \\
\midrule
\textbf{-\bn{র}} & whose & anything \\
\textbf{-\bn{কে}} & the one it happens to & living nouns \\
\textbf{-\bn{দের}} & whose, plural & living nouns \\
\textbf{-\bn{টা}} & the one we mean & things, mostly \\
\bottomrule
\end{tabularx}

All four sit in the same slot at the back of the noun, and the rule from the last chapter still holds: \textbf{one at a time}. That is why \textbf{-\bn{দের}} exists at all, and why you will never see two of these stacked.
\end{grammarlens}

\subsection*{Your turn}
Say these aloud:

\begin{itemize}
\raggedright
  \item \emph{bhāike}, \emph{bônke}, \emph{bôndhuke}
  \item the one that takes no ending, and why — \emph{jāmā}, because it is \textbf{\bn{জিনিস}}
  \item whose, then to whom — \emph{tār bhāi}, then \emph{tār bhāike}
  \item the four endings in the slot — \emph{-r}, \emph{-ke}, \emph{-der}, \emph{-ṭā}
\end{itemize}

\begin{tcolorbox}[breakable,colback=teal!4,colframe=teal!35!black,title={Before you move on}]
What does \textbf{-\bn{কে}} say about a noun? (\textbf{That it is the one the action happens to}.) Which nouns take it? (\textbf{Living ones}.) And how many of these endings can sit on one noun at once? (\textbf{One}.)

Source: \href{https://www.unicode.org/charts/PDF/U0980.pdf}{Unicode Bengali chart}.
\end{tcolorbox}
\section[āmāke]{\bn{আমাকে} (āmāke) — me}

\label{lesson:BN-C25-amake}



\begin{quote}
Put the object ending on \textbf{\bn{ভাই}}. Then say \textbf{I} and \textbf{my}.
\end{quote}

\subsection*{You'll want to know: \bn{আমাকে}}
\begin{quote}
\textbf{\bn{আমাকে}} — \emph{āmāke} — \textbf{me}, \textbf{to me.}
\end{quote}

A pronoun is a living noun, so it takes the living noun's ending — the third time this book has said that and the third time it has been true without an exception:

\noindent
\begin{tabularx}{\linewidth}{@{}>{\raggedright\arraybackslash}X>{\raggedright\arraybackslash}X>{\raggedright\arraybackslash}X@{}}
\toprule
\textbf{job} & \textbf{the word} & \textbf{ending} \\
\midrule
I & \textbf{\bn{আমি}} & — \\
my & \textbf{\bn{আমার}} & \textbf{-\bn{র}} \\
me & \textbf{\bn{আমাকে}} & \textbf{-\bn{কে}} \\
\bottomrule
\end{tabularx}

Same stem \textbf{\bn{আমা}-}, three endings, no memorisation.

\begin{grammarlens}[title={Grammar Lens: the frame that had no filler}]
Chapter nine taught something odd and important: Bengali does not say \emph{I like tea}. It says, near enough, \textbf{tea is good to me} — the liking is not something you do, it is something that happens \textbf{to} you.

That frame needs a \emph{to me}, and the book had none. It has one now:

\begin{quote}
\textbf{\bn{আমাকে} \bn{চা} \bn{ভালো} \bn{লাগে}} — \emph{āmāke chā bhālo lāge} — \textbf{I like tea} (literally, \emph{tea feels good to me}).
\end{quote}

A whole family of Bengali sentences works this way — liking, needing, feeling cold, being hungry — and every one of them was out of reach for want of this word.
\end{grammarlens}

\subsection*{Your turn}
Say these aloud:

\begin{itemize}
\raggedright
  \item the three forms — \emph{āmi}, \emph{āmār}, \emph{āmāke}
  \item \emph{āmāke chā bhālo lāge}
  \item the plural of the same — \emph{āmāder}
  \item on a noun, for comparison — \emph{bhāike}
\end{itemize}

\begin{tcolorbox}[breakable,colback=teal!4,colframe=teal!35!black,title={Before you move on}]
What stem do \textbf{\bn{আমি}}, \textbf{\bn{আমার}} and \textbf{\bn{আমাকে}} share? (\textbf{\bn{আমা}-}.) What does Bengali say instead of \emph{I like tea}? (\textbf{Tea feels good to me}.) And which ending makes the \emph{to me}? (\textbf{-\bn{কে}}.)

Source: \href{https://www.unicode.org/charts/PDF/U0980.pdf}{Unicode Bengali chart}.
\end{tcolorbox}
\section[tomāke]{\bn{তোমাকে} (tomāke) — you, as the one it happens to}

\label{lesson:BN-C25-tomake}



\begin{quote}
Say \textbf{me}. Then say the plain and the careful words for \textbf{you}.
\end{quote}

\subsection*{You'll want to know: \bn{তোমাকে}}
\begin{quote}
\textbf{\bn{তোমাকে}} — \emph{tomāke} — \textbf{you}, as the one it is done to.
\end{quote}

And with the third person, built from the stem you already met in \textbf{\bn{তার}}:

\begin{quote}
\textbf{\bn{তাকে}} — \emph{tāke} — \textbf{him}, \textbf{her.}
\end{quote}

Now the whole table, and it is worth looking at as a table rather than as a list:

\noindent
\begin{tabularx}{\linewidth}{@{}>{\raggedright\arraybackslash}X>{\raggedright\arraybackslash}X>{\raggedright\arraybackslash}X@{}}
\toprule
\textbf{person} & \textbf{possessive} & \textbf{object} \\
\midrule
\textbf{\bn{আমি}} & \textbf{\bn{আমার}} & \textbf{\bn{আমাকে}} \\
\textbf{\bn{তুমি}} & \textbf{\bn{তোমার}} & \textbf{\bn{তোমাকে}} \\
\textbf{\bn{সে}} & \textbf{\bn{তার}} & \textbf{\bn{তাকে}} \\
\bottomrule
\end{tabularx}

Three stems, three endings, nine words, nothing irregular.

\begin{grammarlens}[title={Grammar Lens: what English makes you memorise and Bengali does not}]
The English column is \emph{I, my, me — he, his, him}, and there is no rule in it: a learner memorises each cell. The Bengali column is one stem with \textbf{-\bn{র}} or \textbf{-\bn{কে}} stuck on. Nothing about \textbf{\bn{তাকে}} has to be learned once \textbf{\bn{তার}} is known.

The careful rungs behave, too: \textbf{\bn{আপনাকে}} and \textbf{\bn{তাঁকে}}.

And you can now be polite to a specific person, which chapter twenty-two could not do — it taught \textbf{\bn{দয়া} \bn{করে}} with nobody to address it to.
\end{grammarlens}

\subsection*{Your turn}
Say these aloud:

\begin{itemize}
\raggedright
  \item the object row — \emph{āmāke}, \emph{tomāke}, \emph{tāke}
  \item one stem, two endings — \emph{tār}, then \emph{tāke}
  \item the careful ones — \emph{āpnāke}, \emph{tā̃ke}
  \item politely, to you — \emph{doyā kore} … \emph{tomāke}
\end{itemize}

\begin{tcolorbox}[breakable,colback=teal!4,colframe=teal!35!black,title={Before you move on}]
What is the object form of \textbf{\bn{সে}}? (\textbf{\bn{তাকে}}.) What do \textbf{\bn{তার}} and \textbf{\bn{তাকে}} share? (\textbf{The stem \bn{তা}-}.) And how many of these nine forms are irregular? (\textbf{None}.)

Source: \href{https://www.unicode.org/charts/PDF/U0980.pdf}{Unicode Bengali chart}.
\end{tcolorbox}
\section[ke]{\bn{কে} (ke) — who?}

\label{lesson:BN-C25-ke-who}



\begin{quote}
Put the object ending on \textbf{\bn{ভাই}}. Then say \textbf{you}, as the one it is done to.
\end{quote}

\subsection*{You'll want to know: \bn{কে} standing alone}
\begin{quote}
\textbf{\bn{কে}?} — \emph{ke?} — \textbf{who?}
\end{quote}

\begin{quote}
\textbf{\bn{সে} \bn{কে}?} — \emph{she ke?} — \textbf{who is he?}
\end{quote}

This is a different word from the ending you learned two lessons ago. They are spelled identically and they are not related, and Bengali tells them apart the way it tells \textbf{\bn{ও}} apart from \textbf{\bn{ও}}: by where they sit.

\noindent
\begin{tabularx}{\linewidth}{@{}>{\raggedright\arraybackslash}X>{\raggedright\arraybackslash}X@{}}
\toprule
\textbf{where} & \textbf{what it is} \\
\midrule
stuck to the back of a noun & \textbf{-\bn{কে}}, the object ending \\
standing on its own & \textbf{\bn{কে}}, \emph{who?} \\
\bottomrule
\end{tabularx}

You now have both of the questions a stranger opens with. \textbf{\bn{এটা} \bn{কী}?} asks about a thing; \textbf{\bn{কে}?} asks about a person — and the two are told apart only by their vowel, so say them carefully.

\begin{grammarlens}[title={Grammar Lens: the asking head, for the third time}]
Chapter thirty laid out three heads on one frame — the relative \textbf{\bn{য}-}, the demonstrative \textbf{\bn{ত}-}, and the asking \textbf{\bn{ক}-}. Every question word in this book starts with the last of them:

\begin{itemize}
\raggedright
  \item \textbf{\bn{কী}} — what
  \item \textbf{\bn{কেমন}} — how
  \item \textbf{\bn{কখন}} — when
  \item \textbf{\bn{কে}} — who
\end{itemize}

Four words, one letter at the front of each. And the answers come from the \textbf{\bn{ত}-} column: \textbf{\bn{তিনি}}, \textbf{\bn{তার}}, \textbf{\bn{তখন}}. Ask with \textbf{\bn{ক}-}, answer with \textbf{\bn{ত}-}.
\end{grammarlens}

\subsection*{Your turn}
Say these aloud:

\begin{itemize}
\raggedright
  \item \emph{ke?}
  \item \emph{she ke?} — and answer it — \emph{she āmār bhāi}
  \item the two that sound alike — \emph{bhāike} (the ending), then \emph{ke?} (the question)
  \item the thing and the person — \emph{eṭā ki?}, then \emph{ke?}
  \item the four asking words — \emph{ki}, \emph{kemon}, \emph{kôkhon}, \emph{ke}
\end{itemize}

\begin{tcolorbox}[breakable,colback=teal!4,colframe=teal!35!black,title={Before you move on}]
How do you tell the question \textbf{\bn{কে}} from the ending \textbf{-\bn{কে}}? (\textbf{By where it sits} — alone, or stuck to a noun.) Which letter opens every question word in this book? (\textbf{\bn{ক}-}.) And which column do the answers come from? (\textbf{\bn{ত}-}.)

Source: \href{https://www.unicode.org/charts/PDF/U0980.pdf}{Unicode Bengali chart}.
\end{tcolorbox}
\section[-ke · āmāke · tomāke]{The one it happens to}

\label{lesson:BN-R25-the-one-it-happens-to}



\begin{quote}
Two things in this chapter are written the same and are not the same word. Name both, and say what tells them apart.
\end{quote}

\subsection*{Your turn}
\begin{itemize}
\raggedright
  \item \emph{Say it:} the object row — \emph{āmāke}, \emph{tomāke}, \emph{tāke}
  \item \emph{Say it:} on nouns — \emph{bhāike}, \emph{bônke}, \emph{bôndhuke}
  \item \emph{Say it:} the frame that waited twenty-seven chapters — \emph{āmāke chā bhālo lāge}
  \item \emph{Recall:} two chapters back, the ending that pluralises a living noun — \emph{-rā}
\end{itemize}

\subsection*{Your turn: the distant band — five chapters of grammar, run cold}
Four endings now compete for one slot on the back of a Bengali noun, and one question decides which of them a noun may take. Run all of it.

Say these aloud:

\begin{itemize}
\raggedright
  \item the four endings, and what each says — \emph{-ṭā}, \emph{-r}, \emph{-ke}, \emph{-der}
  \item the question that chooses between them — is it \textbf{\bn{জীব}} or \textbf{\bn{জিনিস}}?
  \item a living noun through all of them — \emph{bhāiṭi}, \emph{bhāir}, \emph{bhāike}, \emph{bhāirā}
  \item a thing, which takes fewer — \emph{jāmāṭā}, \emph{jāmāgulo}
  \item pointing, counting, colouring — \emph{ei ekṭā lāl jāmā}
  \item the six persons and their objects — \emph{āmi/āmāke}, \emph{tumi/tomāke}, \emph{she/tāke}
  \item whose, plural — \emph{āmāder}, \emph{tāder}
  \item the two questions a stranger opens with — \emph{eṭā ki?} and \emph{ke?}
  \item liking, for two people — \emph{āmāke chā bhālo lāge}, \emph{tomāke jôl bhālo lāge}
\end{itemize}

\begin{tcolorbox}[breakable,colback=teal!4,colframe=teal!35!black,title={Before you move on}]
Which nouns take \textbf{-\bn{কে}}? (\textbf{Living ones}.) How many endings may share the slot at once? (\textbf{One}.) What does Bengali say instead of \emph{I like tea}? (\textbf{Tea feels good to me}.) And what is the difference between the two words spelled \textbf{\bn{কে}}? (\textbf{Where they sit} — one is stuck on, one stands alone.)

Source: \href{https://www.unicode.org/charts/PDF/U0980.pdf}{Unicode Bengali chart}.
\end{tcolorbox}
